\subsection*{6.2.}
尚需计算 $\mathcal O_V$ 的上同调群. 置

\begin{equation}\tag{6.2.1}
\mathfrak F_q
=\text{$\mathcal O_V$ 中取 $q$ 次幂的运算}.
\end{equation}

在 $\mathbb P^{n+1}=\mathbb P$ 上有层的正合列,

\begin{equation}\tag{6.2.2}
0\longrightarrow\mathcal O_{\mathbb P}(-d)
\xrightarrow{\,f\,}\mathcal O_{\mathbb P}
\longrightarrow\mathcal O_V\longrightarrow0,
\end{equation}

且下图交换:

\begin{equation}\tag{6.2.3}
\begin{tikzcd}[column sep=large,row sep=large]
0 \arrow[r] &
\mathcal O_{\mathbb P}(-d) \arrow[r,"f"] \arrow[d,"f^{q-1}\mathfrak F_q"'] &
\mathcal O_{\mathbb P} \arrow[r] \arrow[d,"\mathfrak F_q"'] &
\mathcal O_V \arrow[r] \arrow[d,"\mathfrak F_q"'] &
0\\
0 \arrow[r] &
\mathcal O_{\mathbb P}(-d) \arrow[r,"f"] &
\mathcal O_{\mathbb P} \arrow[r] &
\mathcal O_V \arrow[r] &
0 .
\end{tikzcd}
\end{equation}

考虑到

\begin{equation}\tag{6.2.4}
\begin{cases}
H^i(\mathbb P^{n+1},\mathcal O_{\mathbb P})=0,
&i\ne0,\\
H^i(\mathbb P^{n+1},\mathcal O_{\mathbb P}(-d))=0,
&i\ne0,n+1,
\end{cases}
\end{equation}

与 (6.2.2) 相关的上同调长正合列给出

\begin{equation}\tag{6.2.5}
\begin{cases}
H^0(V,\mathcal O_V)\simeq\mathbb F_q,\\
H^i(V,\mathcal O_V)=0,&i\ne0,n,\\
H^n(V,\mathcal O_V)
\simeq H^{n+1}(\mathbb P^{n+1},\mathcal O_{\mathbb P}(-d)).
\end{cases}
\end{equation}

因此交换图 (6.2.3) 给出

\begin{equation}\tag{6.2.6}
\det\!\left(1-t\mathfrak F_q
\mathrel{\Big|}H^0(V,\mathcal O_V)\right)=1-t
\end{equation}

及

\begin{equation}\tag{6.2.7}
\det\!\left(1-t\mathfrak F_q
\mathrel{\Big|}H^n(V,\mathcal O_V)\right)
=
\det\!\left(
1-tf^{q-1}\mathfrak F_q
\mathrel{\Big|}
H^{n+1}(\mathbb P^{n+1},\mathcal O_{\mathbb P}(-d))
\right).
\end{equation}

只需证明

\begin{equation}\tag{6.2.8}
\begin{aligned}
&\det\!\left(
1-tf^{q-1}\mathfrak F_q
\mathrel{\Big|}
H^{n+1}(\mathbb P^{n+1},\mathcal O_{\mathbb P}(-d))
\right)\\
&\qquad=
\det\!\left(
1-t\,\psi_q\circ(X_0f(X))^{q-1}
\mathrel{\Big|}W_S(1)
\right),
\end{aligned}
\end{equation}

这由下述命题推出.

